% Generated draft: requires mathematical and source review.
\documentclass[11pt]{article}
\usepackage[margin=25mm]{geometry}
\usepackage{fontspec}
\setmainfont{Latin Modern Roman}
\usepackage{amsmath,amssymb,mathtools,unicode-math}
\setmathfont{Latin Modern Math}
\usepackage{enumitem,needspace,xurl,hyperref}
\setlength{\parindent}{0pt}
\setlength{\parskip}{5pt}
\setlength{\emergencystretch}{3em}
\newcommand{\N}{\mathbb{N}}
\newcommand{\Z}{\mathbb{Z}}
\newcommand{\Q}{\mathbb{Q}}
\newcommand{\R}{\mathbb{R}}
\newcommand{\C}{\mathbb{C}}
\newcommand{\F}{\mathbb{F}}
\newcommand{\A}{\mathbb{A}}
\newcommand{\PP}{\mathbb{P}}
\newcommand{\E}{\mathbb{E}}
\newcommand{\eps}{\varepsilon}
\newcommand{\dd}{\,\mathrm{d}}
\begin{document}
\section*{UnsolvedMath ID: 20002476}
\textbf{Source problem code:} catalogue entry
\subsection{Definitions and mathematical statement}
Walter Craig's question concerns an area-preserving annulus twist map: for rational rotation $p/q$ there are minimal and minimax periodic Birkhoff orbits; as $p_j/q_j\to\omega\notin\Q$, the source describes convergence of minimal orbits to the Mather set $\Pi_\omega$, and convergence of minimax orbits to the same set when it is an invariant circle. It asks for the fate of minimax orbits when $\Pi_\omega$ is a Cantor set. The topology of convergence and the selection/centering of minimax representatives are not specified.

For the aid's variational reconstruction, suppose the twist map admits a suitable smooth generating function $S(x,y)$, periodic under $(x,y)\mapsto(x+1,y+1)$, with the usual monotone-twist and coercivity hypotheses. A lifted $(p,q)$ configuration satisfies $x_{i+q}=x_i+p$ and has total action
\[
 A_{p,q}(x)=\sum_{i=0}^{q-1}S(x_i,x_{i+1}).
\]
Minimal configurations minimize this action, while minimax configurations are the corresponding mountain-pass critical configurations; the precise minimax class is not given in the source. A Birkhoff configuration is ordered with each of its integer translates $(x_{i+k}+\ell)_i$.

The aid attributes to Mather's classical irrational minimax construction a gap orbit homoclinic to the cantorus: an orbit outside $\Pi_\omega$ whose distance to $\Pi_\omega$ tends to zero in both time directions. This is a gap defect, not an assertion that an additional minimizing Cantor set is produced. Full hypotheses of that construction are not supplied.

Its reported two-scale refinement assumes periods $q_j\to\infty$, rational rotations approaching $\omega$, and total action excess
\[
 \Delta A_j=A_{p_j,q_j}(x^{\mathrm{mm},j})-
 \min A_{p_j,q_j}=o(q_j).
\]
A bounded finite action barrier is a sufficient special case of this sublinear assumption. Under the needed orbit compactness and generating-function hypotheses, the empirical measures
\[
 \mu_j=\frac1{q_j}\sum_{i=0}^{q_j-1}\delta_{f^i(z_j)}
\]
converge weakly to the unique irrational Mather minimizing measure. Nevertheless defect-centered, integer-normalized lifted configurations can have subsequential limits that are stationary Birkhoff configurations outside the minimizing set. Stationarity means
\[
 \partial_2S(x_{i-1},x_i)+\partial_1S(x_i,x_{i+1})=0\quad(i\in\Z).
\]
Measure convergence therefore does not determine every pointed configuration limit.

If the rational minimax critical points are Morse of index one, the aid states that every finite-window, fixed-boundary Hessian of a pointed limit has negative index at most one. If, on one fixed finite window, normalized perturbations retain quadratic form at most $-\delta$ for some $\delta>0$ along the sequence, their limit supplies a negative direction. The pointed limit is then non-minimizing and lies outside the cantorus. No such persistent negative mode is asserted for every sequence; without it a pointed limit may be minimizing.
\subsection{Short English statement}
The source asks what replaces circle convergence of minimax orbits when the irrational Mather set is a cantorus. The aid distinguishes a gap homoclinic defect from the limiting empirical measure: a localized defect can survive after centering even while its relative weight in a long periodic orbit vanishes.
\subsection{Sources}
\begin{itemize}
\item Local UnsolvedMath registry, permanent ID 20002476, source code catalogue entry; created and updated 14 August 2026. Canonical \texttt{aim-pdes-notes.json}, zero-based \texttt{notes[87]}, tag \texttt{problem}, extracted item 4 attributed to Walter Craig.
\item AIM ``New connections between dynamical systems and PDE's'': \url{https://aimath.org/WWN/dynpde/dynpde.pdf}. The supplied verification calls this a damaged extraction of Problem 14, whereas the canonical extract starts with 4 and the background item is 4. This numbering discrepancy is retained. Original remarks and literature field are empty; registry statement status is exact.
\item Research attempt 1: \texttt{partial\_result}, problem status at run \texttt{partially\_solved}, rigor label \texttt{proved}, novelty confidence low. It invokes Mather's irrational minimax construction and finite action barrier without a full publication title, theorem statement, date, or authoritative bibliographic link.
\item Ulam AI / UnsolvedMath Contributors, CC BY 4.0. Machine-generated research aid, not peer reviewed; independently verify before citation.
\end{itemize}
\end{document}
